\section[第8章\q 金钱迷思]{第8章\\ 金钱迷思}
\label{sec:ch8}

搭L线去公司时，柔伊的脑袋嗡嗡作响。

她怎么会觉得亨利是位咖啡师？不是芭芭拉说的？她回想和芭芭拉周一的午餐。不，芭芭拉的原话是：就是你在咖啡厅遇见的老先生啊！他通常早上会在店里煮咖啡。这样而已。是柔伊自己假设亨利是名员工，大概因为看到他时，他总忙着替客人煮咖啡。他看起来确实没有老板架子。

芭芭拉一定知道的，对吗？她怎么也不告诉她一下。她让柔伊去找他谈谈，到底是为了什么呢？柔伊想一道公司就找老板问一问。

可是她并没有。一到33楼之后，她却尽力避着老板。

这天是周四。从周一和芭芭拉一起用午餐之后，她们就没有聊天过，这不是偶然的。也不是柔伊想要逃避，但是想到要讨论离职，去洁西卡介绍的上城区经纪公司工作，她就觉得很不自在。

柔伊自己也没有确定，这就是她还不想面对芭芭拉的原因了。一方面，她必须承认，那边是个不可多得的好机会。她已经厌烦了每天固定的流程，永远都有做不完的事，看不到头。甚至在睡梦中，都还会出现工作内容。而助学贷款每个月还也不见变少，有时候，她甚至觉得还变多了。

这叫她想起乔琪雅的话来，复利是一把双面刃，贷款也是复利的，一旦启动，增长惊人。

哈，柔伊想到，再来说说复利这件事吧！显然乔琪雅和拜伦听了亨利的，至少可说是部分听从，他们的人生已经大有改观了。

对这点，柔伊非常感恩。因为亨利真的好有智能。但是，他的话有些柔伊并不了解，困扰着她。

所以她一边在脑中构想着，一边替春季专辑选图片、修文句。

\subsection{钱赚多，不一定能解决财务问题}

“我们要不要去吃饭啊？”

柔伊转过椅子，看着芭芭拉。已经过下午1：00了吗？

“谢谢，芭芭拉。我今天不出去吃了。”她抽屉里还有一根营养棒，今天应该要吃掉。

“你决定就好。”

柔伊转回她的办公桌。一分钟后，她又回过头来。芭芭拉依然面对着她，朝她这边看着。

“所以，你和亨利谈过了？”

柔伊叹了一口气。“是的。芭芭拉。你一开始为什么不告诉我，他是咖啡厅的老板？”

“那是他的咖啡厅，”芭芭拉说，“又不是我的。再说，你也没问啊！”

“很幽默。”柔伊想了一下，又开口说，“还有，你为什么叫我去找他呢？”

芭芭拉耸肩。“正如我说的，他看事情的方式比较不一样。你总是说，不够钱买这个、买那个的，我都听腻了。”

柔伊笑出来，同时她留意到芭芭拉也快笑了。

“此外，我也没叫你去。或许是，鼓励了一下。”

“鼓励了好大一下。”柔伊回答。

“好吧，好大一下。”芭芭拉顿了一下，“所以呢？”

“所以什么？”

“所以这几天和亨利聊了之后，你有什么想法？”

柔伊叹气。“我还不知道啊，芭芭拉。”她看了看电脑屏幕，又转头过去对着她老板：“在金钱方面，我比较迟钝。”

芭芭拉走到办公室的隔板前方，背靠着隔板，慢慢摇头。“柔伊啊。”

“什么？”柔伊放软了声音。

“这么说吧，我给自己一条规矩，那就是不要介入同事的私人生活。但是我告诉你一句每个女人都该知道的话，好吗？”

“好的，老大。”柔伊说。

芭芭拉不喜欢被称为老大，但是现在她不想离题。“我认真的，你到底要不要听？”

“我洗耳恭听。”柔伊说。

芭芭拉走进柔伊的隔间，靠着办公桌，“我要说的是关于金钱的迷思。这种东西在新闻系不会教，在工商管理学院也不教。”她看了看柔伊的办公桌，“你需不需要记下来啊？”

柔伊打开她的笔电，双手放到打字的位置，看着芭芭拉。“我现在要说的是，学校没有告诉你的金钱迷思。”

芭芭拉点点头。“来，第一条是，赚得多，就会变富有。”

柔伊打开了一个新的文件，用标题字体打出：

{\kaishu 迷思一：赚得多，比较富有}

“亨利可能说过这一则了，对不对？”芭芭拉说，“赚多赚少，和财务稳定没有必然关系。”

“许多人都觉得收入是他们的问题所在，其实不是，\textbf{支出才是问题}。不要误会我的意思，正常收入当然好。但是，\textbf{希望用多赚来解决财务问题，未必行得通}。”

柔伊打字时，尽量保持面无表情。比方说，不露出愧疚之意。芭芭拉会不会知道她在考虑转换跑道的事？很有可能。纽约的媒体圈就像一所高中：每个人都有办法知道其他人的事。

“当你的收入增加，”芭芭拉继续说，“原本的财务问题依然存在，收入增加只是把问题放得更大。因为财务问题关键在于支出习惯，而收入是不会改变花钱习惯的。因此，财务问题的解决方案是养成新习惯，而不是去赚更多钱。”

\subsection{不需要一笔现金，才能理财}

“到这里，你都没问题吧？”

“没问题。”柔伊一边说，一边打字。尽管她并不是真的没问题。当然，这一路下来的思路都很顺，但是她的脑袋一时还转不过来。她的金钱问题，就是她缺钱。为什么多赚了，不能解决问题？

“家庭的理财不是高收入可以解决的，而是如何聪明花钱，如何投资理财。刚好我们就转到了第二个迷思，理财的前提是需要大笔现金。‘钱滚钱’，这话你听过吧？”

柔伊点头。

“这是一个错误的观念。因此我们称之为迷思二。”

柔伊按了分页符号，移到新一页的开始，打下标题：

{\kaishu 迷思二：需要现金，才能钱滚钱}

“我想投资，但是没有钱啊！”你知道，我听见多少女人这么说吗？ ”芭芭拉继续说，“我每次听了都想尖叫。就好像别人说‘若没有意外之财，我哪会发财’是一样的。好像理财是一家昂贵的俱乐部，需要很大一笔入会费似的。

“这些都是错误观念。\textbf{你不需要一笔现金，才能累积财富}。亨利有没有画一张图表给你看？1天存5美元、1天10美元之类的。”

柔伊点头。

“他不是画大饼，那些数字是不会说谎的。复利诚如万有引力。那些创建了财务基础的人，也都是一次存一点，慢慢累积出来的。而且也没风险。只是要下一个决断，面对自己的处境，开始行动。”

“而且你得停止告诉自己，‘我对金钱很不在行。’因为不需要数学天才或华尔街奇才才能理财，只要可以诚实面对自己就好。当然，在当今时代要做到这点不容易。不过，柔伊你可以。这就是我喜欢你的地方，你是诚实的人。”

柔伊觉得自己脸红了。她把双手从键盘上移开，“我懂你说的，芭芭拉。我不想每天的生活都绕着钱转。”

“当然不要。”芭芭拉说，“现在是说，让理财步上轨道，就不会一直为缺钱烦恼。这件事其他人无法代替你。”

\subsection{钱财不会船到桥头自然直}

“这又是一个重点了，柔伊。我们下面的迷思可能是最大的迷思，那就是如果遇到真正的困难，就会出现一个丈夫，或者一位人生导师，或一个英俊潇洒，带着宝剑的骑士，而不是你自己，为自己织一张安全网。其他人会来照顾你。”

柔伊打字。

{\kaishu 迷思三：其他人会来照顾你}

“大多数人也不会直接这么说，”芭芭拉说，“但是他们的行为和选择，却替他们说出来了。‘我男朋友、我先生、我老爸、我的财物顾问，会帮忙我理财’或者是‘船到桥头自然直’。我坦白告诉你，\textbf{金钱这种事，是没有船到桥头自然直的}。”

柔伊想起乔琪雅所说的“我不知道我们透支的情形有多严重”，她也想起了自己妈妈说的话，都是你老爸在管。

“没有带着大袋现金的白马王子。你必须做拯救自己的王子，柔伊。”

“当然对男人来说，这点也是一样。这世上有许多男人也在期待其他人——他们的律师、他们的经纪人、他们的公司、下一任美国总统——来照看他们的财务未来。这样多省事，问题是指望不了别人。”

“你的财富就像你的健康，健康问题不是突然发生的，是要你在生活中每天不断去维护的。你也不可能指望别人来照顾你的健康，财富也一样。他们都完全掌握在你自己手上，任何人无法代劳。”

柔伊打完字，思考了一下。她抬头看着芭芭拉，“可是，你刚刚说这是每个女人都应该知道的。”

芭芭拉点头说：“让我再告诉你一些关于我们的事，关于我们女人。”

“在我们这样开明的年代，女人的收入依然要比男人低20％。公司裁员时，女性员工更易遭资遣，每个女人平均要花10年照顾家里的小孩和老人，因此退休金帐户累积结果要比男人低34％，因此她们也享受较低的社会保险福利。

还有一点：因为女人的平均寿命要比男人长7年，而婚姻有一半是以离婚收场，因此有很大比例的女人晚年是自己度过的。有80％的已婚男子死亡时，尚有婚姻；而有80％已婚女性死亡时已是寡居；80％寡居的女人是在丈夫过世后生活在贫困中。”

芭芭拉停顿了一下，好让柔伊跟上她。

柔伊不知道为什么芭芭拉要说这些，这些真叫人伤感。因为让她想起了老妈，她最近看起来好疲倦，反复发作的背痛，虽然她没说，但是有回老爸在电话里说起了，然后又染上了感冒。如果有一天老爸真的走了，老妈要怎么办？

“你跟得上我说的吗？”芭芭拉问。

柔伊抬起头。“跟上了。可这并不是你说的一个小问题。”

“不是。”芭芭拉说，“而且也不是什么愉快的问题。虽然非常重要。我知道，以你现在二十几岁的年纪，颇难领会其重要性。退休看起来是很老以后的事情。我懂，看起来这些都太远了，但其实人生真的也就是眨眼的功夫。有无数女人经历过有一天突然醒悟过来，发现自己单身一人，一无所有，过去可以有的选择，现在都不存在了，这时她们就会自问：‘怎么回事？我怎么走到今天这一步？’”

柔伊的思绪飞到了周一早上，看到电子屏幕上搁浅的船，她可以想见如果这艘船有个船长，他一定也会问自己同样的问题：“出了什么事，我怎么跑到这里来？”

芭芭拉回去位置后，柔伊低头工作了大约1个小时，突然意识到，她还有一个问题想问亨利，一个如此明显的问题竟然没问，她恨不得戳一戳自己的额头。

今天早上，她问了亨利许多问题，什么10％可不可能啊，关于缴税啊，但那些都是杰弗瑞的问题，不是她的。杰弗瑞的所有疑问中，最大的就是一位70岁的咖啡师，如何知道理财的奥秘？这点，乔琪雅已经说明白了。这位老先生不是咖啡师，而是老板。而他“财务自动化”的讨论，似乎也解答了柔伊关于“如何先付钱给自己”的问题。

但这些都不是现在悬在她脑中的问题。

她需要再找亨利谈一谈。

那天她要回来上班时，咖啡厅的年轻店员工怎么说的？“他想哪时候走就哪时候走。”“通常是3:00，但是也可能更晚一点，也可能更早一点。”

柔伊看了一眼笔电上的时间，下午2:15。如果立即赶过去，说不定还可以堵到他？

柔伊不解自己为什么这么急着见到亨利。早一点晚一点真的很重要吗？反正，她就是必须去问他，好知道答案。她得马上不去找他。

柔伊把笔电收到背包里，冲去找芭芭拉让老板知道她要提早离开。“有点事情。”她一边说着，一边就往电梯冲过去。
